\documentclass[12pt]{article}
\usepackage[utf8]{inputenc}
\usepackage[margin=0.75in]{geometry}
\usepackage{amsmath}
\usepackage{amssymb}
\usepackage{amsthm}
\usepackage{graphicx} % Required for inserting images
\usepackage{array}
\usepackage{tikz-cd}


\title{MAT1100 Homework 1}
\author{Aaron Avram}
\date{September 15 2026}

\begin{document}
\newtheorem{definition}{Definition}
\newtheorem{theorem}{Theorem}
\newtheorem{claim}{Claim}
\newtheorem{corollary}{Corollary}

\maketitle

\begin{align*}
    \textbf{Question 2}
\end{align*}

Let $G \curvearrowright X \neq \varnothing$ and let $G_x$ denote the stabilizer of $x \in X$.
Let $\rho: G \to \operatorname{Sym}(X)$ denote the permutation representation of the action.
\begin{enumerate}
    \item[(a)]
    \begin{proof}
        Suppose $g \in G$ and $x, y \in X$ such that $g \cdot x = y$. Then clearly $x = g^{-1} \cdot y$.
        \begin{enumerate}
            \item[]
            \item[($\subseteq$)] Suppose $h \in G_y$ it suffices to show $h \in gG_xg^{-1}$ and indeed
            \[ g^{-1}hg \cdot x = g^{-1}h \cdot y = g^{-1} \cdot y = x \]
            hence $g^{-1}hg \in G_x$ and therefore $G_y \subseteq gG_xg^{-1}$
            \item[($\supseteq$)] Suppose $h \in G_x$ then
            \[ ghg^{-1} \cdot y = gh \cdot x = g \cdot x = y \]
            hence $ghg^{-1} \in G_y$ and therefore $gG_xg^{-1} \subseteq G_y$.
        \end{enumerate}
    \end{proof}

    \item[(b)]
    \begin{proof}
        First we show that $\ker \rho = \bigcap_{x \in X}G_x$.
        \begin{enumerate}
            \item[]
            \item[($\subseteq$)] If $g \in \ker \rho$ then for all $x \in X$, $g \cdot x = x$.
            Hence $g \in G_x$ for all $x \in X$ and so
            \[ g \in \bigcap_{x \in X}G_x \]
            \item[($\supseteq$)] If $g \in \bigcap_{x \in X}G_x$ then for all $x \in X$
            $ g \in G_x$ and therefore $g \cdot x = x$ for all $x \in X$. In other words, $g \in \ker \rho$.
        \end{enumerate}
        If $G \curvearrowright X$ is transitive then given $x \in X$
        \[ \mathcal O_x = \{ g \cdot x : g \in G \} = X \]
        and by $(a)$, $\mathcal O_{gx} = g\mathcal O_x g^{-1}$. Thus
        \[ \ker \rho = \bigcap_{x \in X}G_x = \bigcap_{g \in G}G_{g \cdot x} = \bigcap_{g \in G}gG_xg^{-1} \]
    \end{proof}

    \item[(c)]
    \begin{proof}
        Fix $x \in X$. Let $\phi: G \to X$ be given by
        \[ \phi(g) = g \cdot x \]
        First we show that for all $y \in X$ there exists unique $g \in G$ such that
        \[ g \cdot x = y \]
        the existence of such a $G$ is evident as $G \curvearrowright X$ is transitive.
        Therefore $\phi$ is surjective.

        Now for uniqueness, suppose $g, h \in G$ such that
        \[ g \cdot x = h \cdot x \]
        then
        \[ g^{-1}h \cdot x = x \]
        hence $g^{-1}h \in G_x$. Now given $a \in G_x$, $gag^{-1} = a$
        as $G$ is abelian hence $gG_xg^{-1} = G_x$ for all $g \in G$. Since $G \curvearrowright X$ is transitive
        \[ \ker \rho = \bigcup_{g \in G} gG_xg^{-1} = \bigcap_{g \in G}G_x = G_x \]
        hence $g^{-1}h \in G_x = \ker \rho$ and since $\rho$ is faithful this implies that
        $g^{-1}h = e$ and $g = h$.

        The uniqueness of the above element $g$ in turn implies that $\phi$ is injective.
        Therefore $\phi$ is a bijection.
    \end{proof}
\end{enumerate}

\end{document}